% Options for packages loaded elsewhere
\PassOptionsToPackage{unicode}{hyperref}
\PassOptionsToPackage{hyphens}{url}
\documentclass[
  nobib]{tufte-handout}
\usepackage{xcolor}
\usepackage{amsmath,amssymb}
\setcounter{secnumdepth}{-\maxdimen} % remove section numbering
\usepackage{iftex}
\ifPDFTeX
  \usepackage[T1]{fontenc}
  \usepackage[utf8]{inputenc}
  \usepackage{textcomp} % provide euro and other symbols
\else % if luatex or xetex
  \usepackage{unicode-math} % this also loads fontspec
  \defaultfontfeatures{Scale=MatchLowercase}
  \defaultfontfeatures[\rmfamily]{Ligatures=TeX,Scale=1}
\fi
\usepackage{lmodern}
\ifPDFTeX\else
  % xetex/luatex font selection
\fi
% Use upquote if available, for straight quotes in verbatim environments
\IfFileExists{upquote.sty}{\usepackage{upquote}}{}
\IfFileExists{microtype.sty}{% use microtype if available
  \usepackage[]{microtype}
  \UseMicrotypeSet[protrusion]{basicmath} % disable protrusion for tt fonts
}{}
\makeatletter
\@ifundefined{KOMAClassName}{% if non-KOMA class
  \IfFileExists{parskip.sty}{%
    \usepackage{parskip}
  }{% else
    \setlength{\parindent}{0pt}
    \setlength{\parskip}{6pt plus 2pt minus 1pt}}
}{% if KOMA class
  \KOMAoptions{parskip=half}}
\makeatother
\usepackage{color}
\usepackage{fancyvrb}
\newcommand{\VerbBar}{|}
\newcommand{\VERB}{\Verb[commandchars=\\\{\}]}
\DefineVerbatimEnvironment{Highlighting}{Verbatim}{commandchars=\\\{\}}
% Add ',fontsize=\small' for more characters per line
\newenvironment{Shaded}{}{}
\newcommand{\AlertTok}[1]{\textcolor[rgb]{1.00,0.00,0.00}{\textbf{#1}}}
\newcommand{\AnnotationTok}[1]{\textcolor[rgb]{0.38,0.63,0.69}{\textbf{\textit{#1}}}}
\newcommand{\AttributeTok}[1]{\textcolor[rgb]{0.49,0.56,0.16}{#1}}
\newcommand{\BaseNTok}[1]{\textcolor[rgb]{0.25,0.63,0.44}{#1}}
\newcommand{\BuiltInTok}[1]{\textcolor[rgb]{0.00,0.50,0.00}{#1}}
\newcommand{\CharTok}[1]{\textcolor[rgb]{0.25,0.44,0.63}{#1}}
\newcommand{\CommentTok}[1]{\textcolor[rgb]{0.38,0.63,0.69}{\textit{#1}}}
\newcommand{\CommentVarTok}[1]{\textcolor[rgb]{0.38,0.63,0.69}{\textbf{\textit{#1}}}}
\newcommand{\ConstantTok}[1]{\textcolor[rgb]{0.53,0.00,0.00}{#1}}
\newcommand{\ControlFlowTok}[1]{\textcolor[rgb]{0.00,0.44,0.13}{\textbf{#1}}}
\newcommand{\DataTypeTok}[1]{\textcolor[rgb]{0.56,0.13,0.00}{#1}}
\newcommand{\DecValTok}[1]{\textcolor[rgb]{0.25,0.63,0.44}{#1}}
\newcommand{\DocumentationTok}[1]{\textcolor[rgb]{0.73,0.13,0.13}{\textit{#1}}}
\newcommand{\ErrorTok}[1]{\textcolor[rgb]{1.00,0.00,0.00}{\textbf{#1}}}
\newcommand{\ExtensionTok}[1]{#1}
\newcommand{\FloatTok}[1]{\textcolor[rgb]{0.25,0.63,0.44}{#1}}
\newcommand{\FunctionTok}[1]{\textcolor[rgb]{0.02,0.16,0.49}{#1}}
\newcommand{\ImportTok}[1]{\textcolor[rgb]{0.00,0.50,0.00}{\textbf{#1}}}
\newcommand{\InformationTok}[1]{\textcolor[rgb]{0.38,0.63,0.69}{\textbf{\textit{#1}}}}
\newcommand{\KeywordTok}[1]{\textcolor[rgb]{0.00,0.44,0.13}{\textbf{#1}}}
\newcommand{\NormalTok}[1]{#1}
\newcommand{\OperatorTok}[1]{\textcolor[rgb]{0.40,0.40,0.40}{#1}}
\newcommand{\OtherTok}[1]{\textcolor[rgb]{0.00,0.44,0.13}{#1}}
\newcommand{\PreprocessorTok}[1]{\textcolor[rgb]{0.74,0.48,0.00}{#1}}
\newcommand{\RegionMarkerTok}[1]{#1}
\newcommand{\SpecialCharTok}[1]{\textcolor[rgb]{0.25,0.44,0.63}{#1}}
\newcommand{\SpecialStringTok}[1]{\textcolor[rgb]{0.73,0.40,0.53}{#1}}
\newcommand{\StringTok}[1]{\textcolor[rgb]{0.25,0.44,0.63}{#1}}
\newcommand{\VariableTok}[1]{\textcolor[rgb]{0.10,0.09,0.49}{#1}}
\newcommand{\VerbatimStringTok}[1]{\textcolor[rgb]{0.25,0.44,0.63}{#1}}
\newcommand{\WarningTok}[1]{\textcolor[rgb]{0.38,0.63,0.69}{\textbf{\textit{#1}}}}
\usepackage{longtable,booktabs,array}
\newcounter{none} % for unnumbered tables
\usepackage{calc} % for calculating minipage widths
% Correct order of tables after \paragraph or \subparagraph
\usepackage{etoolbox}
\makeatletter
\patchcmd\longtable{\par}{\if@noskipsec\mbox{}\fi\par}{}{}
\makeatother
% Allow footnotes in longtable head/foot
\IfFileExists{footnotehyper.sty}{\usepackage{footnotehyper}}{\usepackage{footnote}}
\makesavenoteenv{longtable}
\setlength{\emergencystretch}{3em} % prevent overfull lines
\providecommand{\tightlist}{%
  \setlength{\itemsep}{0pt}\setlength{\parskip}{0pt}}
% =====================================================================
% axona-doc-preamble.tex — shared LaTeX preamble for the programmer-guide
% quartet (Quick-Start, Programmer Guide, API Reference, Services Guide).
%
% These docs are markdown-native; the PDF is built by converting the .md to
% a tufte-handout .tex with pandoc (-H this file) and compiling with tectonic.
% See render.sh. Modeled on explainer/Axona-Explainer.tex (the flat-essay
% original); the tweaks below make tufte-handout work for code/table-heavy
% reference material.
% =====================================================================

\definecolor{rust}{RGB}{185,78,54}
\definecolor{rustsoft}{RGB}{212,168,150}
\definecolor{inkmuted}{RGB}{102,102,102}
\usepackage{microtype}

% Explainer-matched page furniture: tufte-handout's letterspaced running
% heads fail on certain page breaks (the SOUL \MakeTextLowercase issue the
% explainer works around), so use the same fancyhdr style — no running
% head, a centered page number in the footer.
\usepackage{fancyhdr}
\fancypagestyle{ndhtplain}{%
  \fancyhf{}%
  \renewcommand{\headrulewidth}{0pt}%
  \fancyfoot[C]{\thepage}%
}

% Explainer-matched heading-orphan protection: reserve vertical room before
% each (sub)section heading so a heading never sits alone at a page bottom.
\usepackage{etoolbox}
\usepackage{needspace}
\pretocmd{\section}{\needspace{6\baselineskip}}{}{}
\pretocmd{\subsection}{\needspace{4\baselineskip}}{}{}

% XeTeX (tectonic) unicode monospace: renders … → and the box-drawing ASCII
% diagrams that the default CM/LM mono cannot. Menlo ships with macOS.
\setmonofont{Menlo}[Scale=0.78]

% tufte-handout ships \subsubsection as an error stub (it wants flat essays).
% A run-in fallback lets nested reference headings (§4.2, §12.3, …) compile.
\renewcommand{\subsubsection}[1]{\medskip\par\noindent{\normalsize\bfseries #1}\par\nobreak\smallskip}

% These docs use no margin notes, so reclaim tufte's wide right margin: widen
% the text block so code blocks and API tables don't overflow it.
\usepackage{geometry}
\geometry{left=1in, textwidth=6.1in, marginparsep=0.2in, marginparwidth=0.9in,
          top=1in, headsep=0.4in}

% Wrap long code lines (cheat-sheet comments run to ~130 chars) with a ↪ marker
% instead of overflowing. Applies to pandoc's fancyvrb Highlighting/Verbatim.
\usepackage{fvextra}
\fvset{breaklines=true, breakanywhere=true}

% A few technical symbols appear in prose / inline code that Latin Modern and
% Menlo don't cover; map them to math equivalents so they render (rather than
% dropping to tofu). Only maps chars actually used — leaves → alone (it renders).
\usepackage{newunicodechar}
\newunicodechar{≈}{\ensuremath{\approx}}
\newunicodechar{≥}{\ensuremath{\geq}}
\newunicodechar{≤}{\ensuremath{\leq}}
\newunicodechar{×}{\ensuremath{\times}}
\newunicodechar{·}{\ensuremath{\cdot}}
\newunicodechar{‖}{\ensuremath{\|}}
\usepackage{bookmark}
\IfFileExists{xurl.sty}{\usepackage{xurl}}{} % add URL line breaks if available
\urlstyle{same}
\hypersetup{
  hidelinks,
  pdfcreator={LaTeX via pandoc}}

\author{}
\date{}

\begin{document}

% Generated by render.sh — the explainer-style title page.
\thispagestyle{empty}
\null\vfill
\begin{center}
{\fontsize{48}{52}\selectfont \textsc{Axona}\par}
\vspace{0.6em}
{\Large\itshape Bridges, Relays, and the Programs You Run\par}
\vspace{4em}
{\large David A.~Smith\par}
{\itshape\small Axona.net\par}
\vspace{2em}
{\small The Axona Services Guide \quad v4.22.0 \quad 2026-07-14\par}
\end{center}
\vfill
\null
\newpage
\pagestyle{ndhtplain}
\thispagestyle{ndhtplain}
\setcounter{page}{1}

\section{Axona Services Guide}\label{axona-services-guide}

\emph{(kernel 4.22.0 · testnet --- versioned with the protocol it
describes)}

The other programmer-guide documents teach the \textbf{library} --- how
to build your own peer with \texttt{@axona/protocol} and speak pub/sub.
This guide covers the \textbf{services}: the ready-made programs you
\emph{run} (or rent) rather than write --- the signaling
\textbf{bridge}, the \textbf{relay} and its four front-ends (console,
CLI, MCP server, desktop app), the \textbf{directory/federation} that
ties bridges together, and the \textbf{PoW collector}. If you are
operating an Axona deployment, wiring an agent into the network, or just
want a topic to stay alive when no browser tab is open, this is the
document for you.

\begin{itemize}
\tightlist
\item
  \textbf{Protocol kernel}:
  \href{https://github.com/axona-net/axona-protocol}{@axona/protocol}
  (v4.22.0)
\item
  \textbf{Wire version}: 4.0 (\texttt{WIRE\_VERSION}); kernel version
  4.22.0 (\texttt{KERNEL\_VERSION})
\item
  \textbf{Live networks}: both run the 4.x line. \textbf{Testnet}
  (\texttt{wss://testnet.axona.net}) tracks the newest kernel ---
  4.22.0, the version this guide describes. The \textbf{production}
  federated pair (\texttt{wss://bridge.axona.net} east +
  \texttt{wss://bridge-west.axona.net} west) runs the most recently
  promoted kernel (4.21.0 at press time), typically one release behind
  while changes soak on testnet. The two networks are wire-compatible
  but \emph{separate} --- a peer joins one or the other; pick with
  \texttt{RELAY\_NETWORK}/\texttt{BRIDGE\_URL} (§3).
\item
  \textbf{Companion docs}:

  \begin{itemize}
  \tightlist
  \item
    \href{Quick-Start-v4.22.0.md}{Quick Start} ·
    \href{Axona-Programmer-Guide-v4.22.0.md}{Programmer Guide} ·
    \href{Axona-API-Reference-v4.22.0.md}{API Reference} --- the
    library.
  \item
    \href{Axona-AI-Grounding-v4.22.0.md}{AI Grounding} --- building with
    an AI assistant.
  \item
    \href{../architecture/Axona-Architecture.tex}{Architecture} --- how
    the bridge + transport work under the hood.
  \end{itemize}
\end{itemize}

\begin{quote}
\textbf{Library vs.~services.} A \emph{peer} is the unit the library
gives you: an identity + a transport + a pub/sub manager, running inside
a browser tab, a Node process, or an agent. The library is enough to
build a complete app. The \textbf{services} exist because some jobs
outlive a single tab: brokering the first WebRTC handshake between two
strangers (the \textbf{bridge}), keeping a topic's messages available
when every author has closed their laptop (a \textbf{relay} that
\texttt{host()}s the topic), and letting non-browser callers --- shells,
agents, GUIs --- publish and subscribe (the \textbf{CLI / MCP / desktop}
front-ends). None of them is privileged: every service is just a peer
that has volunteered to stay online and take on a role.
\end{quote}

\begin{center}\rule{0.5\linewidth}{0.5pt}\end{center}

\subsection{1. The service map}\label{the-service-map}

\begin{verbatim}
                      +----------------------------------------------+
   browser apps  ---> |  BRIDGE   wss://bridge.axona.net (+ west)     |
   (axona.net,        |  - WebSocket rendezvous (first handshake)     |
    demo apps)        |  - STUN/TURN (NAT traversal, coturn)          |
                      |  - embedded peer: federation uplink + hosts   |
   relays / CLI / --->|    the bridge DIRECTORY topic                 |
   MCP / desktop      +-----------------------+----------------------+
                                              |  (after signaling, peers
                                              v   talk WebRTC; bridge optional)
                      +----------------------------------------------+
                      |  RELAY   "a subset of the bridge"            |
                      |  a headless peer that joins the mesh to:     |
                      |  - host()  keyspace / topics  -> durable roots|
                      |  - publish METRIC snapshots   -> metricTopic  |
                      |  front-ends over one core (src/ops.js):       |
                      |  console/TUI . CLI . MCP server . desktop app |
                      +----------------------------------------------+
\end{verbatim}

{\def\LTcaptype{none} % do not increment counter
\begin{longtable}[]{@{}
  >{\raggedright\arraybackslash}p{(\linewidth - 6\tabcolsep) * \real{0.2500}}
  >{\raggedright\arraybackslash}p{(\linewidth - 6\tabcolsep) * \real{0.2500}}
  >{\raggedright\arraybackslash}p{(\linewidth - 6\tabcolsep) * \real{0.2500}}
  >{\raggedright\arraybackslash}p{(\linewidth - 6\tabcolsep) * \real{0.2500}}@{}}
\toprule\noalign{}
\begin{minipage}[b]{\linewidth}\raggedright
Service
\end{minipage} & \begin{minipage}[b]{\linewidth}\raggedright
Package / entry
\end{minipage} & \begin{minipage}[b]{\linewidth}\raggedright
What it is
\end{minipage} & \begin{minipage}[b]{\linewidth}\raggedright
Runs where
\end{minipage} \\
\midrule\noalign{}
\endhead
\bottomrule\noalign{}
\endlastfoot
\textbf{Bridge} & \texttt{axona-bridge} (\texttt{src/server.js}) &
Signaling rendezvous + STUN/TURN + a federating, directory-hosting peer
& A server (Docker stack) \\
\textbf{Relay --- console/TUI} & \texttt{axona-relay}
(\texttt{src/index.js}, bin \texttt{axona-relay}) & Headless node that
hosts topics + publishes metrics; live dashboard & A box you keep
online \\
\textbf{Relay --- CLI} & \texttt{axona-cli} (\texttt{src/cli.js}) &
One-shot \texttt{pub} / \texttt{sub} / \texttt{pull} from a shell or
script & Anywhere with Node \\
\textbf{Relay --- MCP server} & \texttt{axona-mcp} (\texttt{src/mcp.js})
& A \textbf{persistent peer} with a durable identity, exposed as native
\textbf{agent tools} (publish/subscribe/host) over MCP/stdio & Spawned
by the agent host \\
\textbf{Relay --- desktop app} & \texttt{axona-relay/desktop} (Electron)
& A GUI wrapper around the relay & A laptop/desktop \\
\textbf{PoW collector} & \texttt{axona-relay/pow-collector.js}
(\texttt{npm\ run\ collect}) & Subscribes to the PoW-bench results topic
and archives them & A box you keep online \\
\textbf{Directory / federation} & kernel \texttt{bridgeDirectory.js} & A
\emph{behavior}, not a process: bridges advertise, clients fail over &
Inside bridges + clients \\
\end{longtable}
}

The same \texttt{@axona/protocol} kernel runs in every box above. There
is no special ``server build'' --- a bridge is a peer that also
terminates WebSockets and TURN; a relay is a peer that also
\texttt{host()}s and publishes metrics.

\begin{center}\rule{0.5\linewidth}{0.5pt}\end{center}

\subsection{2. The Bridge}\label{the-bridge}

\subsubsection{What it is}\label{what-it-is}

A \textbf{bridge} is the only piece of fixed infrastructure Axona needs,
and even it is needed only at the \textbf{start} of a connection. Two
browsers that have never met cannot exchange WebRTC offer/answer SDP on
their own --- they need a mutually reachable rendezvous. The bridge is
that rendezvous. It does three jobs:

\begin{enumerate}
\def\labelenumi{\arabic{enumi}.}
\tightlist
\item
  \textbf{WebSocket signaling} --- relays the offer/answer/ICE handshake
  between two peers until their direct WebRTC data channel is open.
  After that the bridge is \textbf{out of the data path}; peers talk
  peer-to-peer over the mesh.
\item
  \textbf{STUN/TURN} --- a bundled \textbf{coturn} gives peers their
  public reflexive address (STUN) and, for the \textasciitilde10--15\%
  of NAT pairs that can't connect directly, relays the media (TURN). The
  bridge mints short-lived TURN credentials in its welcome frame.
\item
  \textbf{An embedded peer} --- the bridge process also \emph{joins its
  own network} as a node. That embedded peer \textbf{federates} (an
  outbound uplink to a sibling bridge, §5) and \textbf{hosts the bridge
  directory topic} so clients can discover other bridges and fail over.
\end{enumerate}

\begin{quote}
A bridge is a convenience, not an authority. It cannot read message
payloads (they're end-to-end between mesh peers), cannot forge a publish
(every envelope is author-signed), and any number of independent bridges
can coexist --- clients rank and fail over between them.
\end{quote}

\subsubsection{Running one}\label{running-one}

Production bridges run as a \textbf{Docker Compose stack} ---
\texttt{bridge} + \textbf{Caddy} (automatic TLS) + \textbf{coturn}
(TURN) --- from the \texttt{axona-bridge} repo. The one-command path for
a fresh Ubuntu/Debian host:

\begin{Shaded}
\begin{Highlighting}[]
\ExtensionTok{curl} \AttributeTok{{-}fsSL}\NormalTok{ https://raw.githubusercontent.com/axona{-}net/axona{-}bridge/main/deploy/install.sh }\DataTypeTok{\textbackslash{}}
  \KeywordTok{|} \FunctionTok{sudo}\NormalTok{ BRIDGE\_DOMAIN=bridge.example.net LETSENCRYPT\_EMAIL=you@example.net bash}
\end{Highlighting}
\end{Shaded}

or, with the repo checked out:

\begin{Shaded}
\begin{Highlighting}[]
\BuiltInTok{cd}\NormalTok{ /opt/axona{-}bridge}
\ExtensionTok{docker}\NormalTok{ compose build        }\CommentTok{\# builds first; the old container keeps serving}
\ExtensionTok{docker}\NormalTok{ compose up }\AttributeTok{{-}d}        \CommentTok{\# fast swap; Caddy keeps its certificate}
\ExtensionTok{curl}\NormalTok{ https://bridge.example.net/healthz     }\CommentTok{\# → \{"status":"ok","version":…,"kernelVersion":"4.22.0"\}}
\end{Highlighting}
\end{Shaded}

Full provisioning options (the installer, the Docker bundle, and a
manual walkthrough) are in
\textbf{\texttt{axona-bridge/deploy/INSTALL.md}}.

\subsubsection{Configuration}\label{configuration}

Bridge config lives in \texttt{/etc/axona-bridge.env} (systemd) or the
compose \texttt{.env}:

{\def\LTcaptype{none} % do not increment counter
\begin{longtable}[]{@{}
  >{\raggedright\arraybackslash}p{(\linewidth - 2\tabcolsep) * \real{0.5000}}
  >{\raggedright\arraybackslash}p{(\linewidth - 2\tabcolsep) * \real{0.5000}}@{}}
\toprule\noalign{}
\begin{minipage}[b]{\linewidth}\raggedright
Var
\end{minipage} & \begin{minipage}[b]{\linewidth}\raggedright
Meaning
\end{minipage} \\
\midrule\noalign{}
\endhead
\bottomrule\noalign{}
\endlastfoot
\texttt{DOMAIN} / \texttt{PUBLIC\_IP} & Public hostname + IP (TLS + TURN
advertisement) \\
\texttt{TURN\_AUTH\_SECRET} & Shared secret with coturn's
\texttt{static-auth-secret}; the bridge mints creds from it \\
\texttt{BRIDGE\_PUBLIC\_URL} & The \texttt{wss://} URL this bridge
advertises in the directory
(e.g.~\texttt{wss://bridge-west.axona.net}) \\
\texttt{BRIDGE\_DIRECTORY} & \texttt{on} (advertise + host the
directory) or \texttt{off} (isolated fleet --- testnet uses
\texttt{off}) \\
\texttt{BRIDGE\_LAT} / \texttt{BRIDGE\_LNG} / \texttt{REGION\_LABEL} &
The bridge's geo anchor + a human label for its directory entry \\
\texttt{MIN\_PEER\_VERSION} & Lowest wire version admitted; older peers
get \texttt{UPGRADE\_REQUIRED} (WS close 4426) \\
\texttt{HOST} & Bind address --- production sets \texttt{127.0.0.1} so
the reverse proxy is the only ingress \\
\texttt{HEALTHZ\_TOKEN} & Gates the \textbf{full} \texttt{/healthz} body
+ \texttt{/diag}; unauthenticated \texttt{/healthz} returns only
\texttt{\{status,version,kernelVersion\}} \\
\end{longtable}
}

\subsubsection{Health \& introspection}\label{health-introspection}

\begin{itemize}
\tightlist
\item
  \texttt{GET\ /healthz} --- unauthenticated, returns
  \texttt{\{status,\ version,\ kernelVersion\}}.
\item
  \texttt{GET\ /healthz} with header
  \texttt{X-Healthz-Token:\ \textless{}token\textgreater{}} --- the full
  body: synaptome size, directory \texttt{\{enabled,url\}}, uplink
  status, peer counts.
\item
  \texttt{GET\ /diag} (token-gated) --- deeper diagnostics.
\end{itemize}

\begin{center}\rule{0.5\linewidth}{0.5pt}\end{center}

\subsection{3. The Relay}\label{the-relay}

\subsubsection{What it is}\label{what-it-is-1}

A \textbf{relay is a subset of the bridge}: the embedded-peer half,
without the WebSocket/TURN server. It is a headless node you keep online
so the network has durable, well-placed members. Two jobs:

\begin{enumerate}
\def\labelenumi{\arabic{enumi}.}
\tightlist
\item
  \textbf{Hosting} --- \texttt{host()} makes the relay a \textbf{root}
  for a keyspace region (or for named topics) \emph{without
  subscribing}, so it stores-and-serves a topic's messages even when it
  doesn't consume them. This is what keeps a topic's backlog answerable
  after every author closes their tab. By default a relay hosts its
  region's keyspace (\texttt{0x89} for \texttt{useast}). Under the 4.x
  \textbf{cohort} model a topic is replicated across the K-closest
  nodes, so durability no longer rests on one relay --- several
  co-hosting relays converge on the same history via anti-entropy, and
  any of them can answer a read.
\item
  \textbf{Metrics} --- served \textbf{on demand} (v4.12.0): when any
  client subscribes to a topic's derived \texttt{metricTopic(T)}, a
  renewable \emph{metrics-on} lease routes to the topic's root --- often
  a relay, since relays root what they host --- and while the lease is
  fresh that root publishes a signed snapshot
  (\texttt{\{current\_count,\ seq,\ subscribers,\ bytes\}}) every
  \textasciitilde20 s, giving subscribers live counts + a rolling
  \textasciitilde48 h trend instead of polling. The first snapshot lands
  \textasciitilde2--20 s after demand turns on; publishing stops
  \textasciitilde70 s after the last metric subscriber leaves. (The
  pre-4.12 push loop --- publish every rooted topic on a timer, watched
  or not --- is retired.) \textbf{Both open and owned} topics serve
  metrics (owned-topic activity counts are public too); \texttt{seq} is
  the dense message counter (total events ever, kills included). Because
  every co-hosting root publishes its own snapshot, the reader-side
  \texttt{peer.metrics()} aggregates across the cohort (sums
  subscribers, maxes the rest --- see the API Reference).
\end{enumerate}

\subsubsection{Running one}\label{running-one-1}

\begin{Shaded}
\begin{Highlighting}[]
\BuiltInTok{cd}\NormalTok{ axona{-}relay}
\ExtensionTok{npm}\NormalTok{ install}
\ExtensionTok{npm}\NormalTok{ start          }\CommentTok{\# live dashboard (interactive TTY)}
\ExtensionTok{npm}\NormalTok{ run probe      }\CommentTok{\# plain timestamped log lines (RELAY\_TUI=0) — for piping / services}
\end{Highlighting}
\end{Shaded}

Run it again in another terminal for a second node --- the
\textbf{first} instance claims the persistent identity (stable nodeId,
\texttt{identity.\textless{}region\textgreater{}.json}); each
\textbf{additional} instance mints a fresh ephemeral identity in the
same region. So you get one well-known node plus as many throwaway nodes
as you want.

\subsubsection{Configuration (env)}\label{configuration-env}

{\def\LTcaptype{none} % do not increment counter
\begin{longtable}[]{@{}
  >{\raggedright\arraybackslash}p{(\linewidth - 4\tabcolsep) * \real{0.3333}}
  >{\raggedright\arraybackslash}p{(\linewidth - 4\tabcolsep) * \real{0.3333}}
  >{\raggedright\arraybackslash}p{(\linewidth - 4\tabcolsep) * \real{0.3333}}@{}}
\toprule\noalign{}
\begin{minipage}[b]{\linewidth}\raggedright
Var
\end{minipage} & \begin{minipage}[b]{\linewidth}\raggedright
Default
\end{minipage} & \begin{minipage}[b]{\linewidth}\raggedright
Meaning
\end{minipage} \\
\midrule\noalign{}
\endhead
\bottomrule\noalign{}
\endlastfoot
\texttt{RELAY\_NETWORK} & \texttt{prod} & Bootstrap network:
\texttt{prod} (\texttt{bridge.axona.net}) or \texttt{testnet}
(\texttt{testnet.axona.net}). Both run the 4.x line; testnet tracks the
newest kernel. Use \texttt{testnet} to follow the staging line this
guide describes; \texttt{prod} for production service. \\
\texttt{BRIDGE\_URL} & --- & Explicit bridge URL; \textbf{overrides}
\texttt{RELAY\_NETWORK} \\
\texttt{RELAY\_REGION} & --- & \texttt{auto} (detect), a region name
(\texttt{useast}), or code (\texttt{0x89}) --- sets the nodeId geo
prefix \\
\texttt{RELAY\_LAT} / \texttt{RELAY\_LNG} &
\texttt{37.77}/\texttt{-122.42} & Geo prefix by coordinate (if
\texttt{RELAY\_REGION} unset). Default = SF (\texttt{uswest}) \\
\texttt{RELAY\_IDENTITY\_PATH} &
\texttt{./identity.\textless{}region\textgreater{}.json} & Persisted
keypair (stable nodeId) \\
\texttt{RELAY\_HOST\_KEYSPACE} & \texttt{1} & Host the region's whole
keyspace (set \texttt{0} to host only \texttt{RELAY\_TOPICS}) \\
\texttt{RELAY\_TOPICS} & --- & Comma-separated topic names to host
explicitly \\
\texttt{RELAY\_METRICS} & \texttt{1} & Publish metric snapshots
(\texttt{0} to disable) \\
\texttt{RELAY\_METRICS\_INTERVAL\_MS} & \textasciitilde20000 & Metric
publish cadence (\textasciitilde20 s) \\
\texttt{RELAY\_TUI} & auto & \texttt{1} force dashboard, \texttt{0}
force plain log \\
\end{longtable}
}

Bridge precedence: \texttt{BRIDGE\_URL} › \texttt{RELAY\_NETWORK} ›
\texttt{prod}. Region precedence: \texttt{RELAY\_REGION} ›
\texttt{RELAY\_LAT}/\texttt{RELAY\_LNG} › default SF.

\subsubsection{Run as a service
(systemd)}\label{run-as-a-service-systemd}

\begin{Shaded}
\begin{Highlighting}[]
\CommentTok{\# /etc/systemd/system/axona{-}relay.service}
\KeywordTok{[Service]}
\DataTypeTok{Environment}\OtherTok{=}\StringTok{RELAY\_REGION=useast RELAY\_TUI=0}
\DataTypeTok{ExecStart}\OtherTok{=}\StringTok{/usr/bin/node /opt/axona{-}relay/src/index.js}
\DataTypeTok{Restart}\OtherTok{=}\StringTok{always}
\end{Highlighting}
\end{Shaded}

\begin{quote}
\textbf{Operator note --- process count.} Each relay forks \textbf{one
worker child}, so a healthy \emph{N}-node fleet is \textbf{2 × N}
\texttt{node\ src/index.js} processes (a parent + worker per node).
Judge fleet health by the distinct \textbf{bridge connections} in the
logs, not the raw process count. If the \textbf{desktop app} (below) is
also running, it carries \emph{its own} embedded relay that respawns if
killed --- expect one extra.
\end{quote}

\begin{center}\rule{0.5\linewidth}{0.5pt}\end{center}

\subsection{4. Relay front-ends}\label{relay-front-ends}

The relay's core pub/sub lives in \texttt{src/ops.js} (\texttt{publish}
/ \texttt{subscribe} / \texttt{pull}). The \textbf{CLI} uses it directly
--- one throwaway peer per call. The \textbf{MCP server} wraps it in a
\emph{persistent} peer (\texttt{src/mcp-session.js}) so an agent can be
a standing participant rather than a series of one-shot connections. The
\textbf{desktop app} is the GUI relay.

\subsubsection{\texorpdfstring{4a. CLI
(\texttt{axona-cli})}{4a. CLI (axona-cli)}}\label{a.-cli-axona-cli}

One-shot pub/sub/pull from a shell or script --- no long-lived process:

\begin{Shaded}
\begin{Highlighting}[]
\ExtensionTok{npm}\NormalTok{ run pub  }\AttributeTok{{-}{-}} \AttributeTok{{-}{-}topic}\NormalTok{ us{-}east/hello }\AttributeTok{{-}{-}message} \StringTok{"hi"}          \CommentTok{\# prints \{ ok, msgId, … \}}
\ExtensionTok{npm}\NormalTok{ run sub  }\AttributeTok{{-}{-}} \AttributeTok{{-}{-}topic}\NormalTok{ us{-}east/hello }\AttributeTok{{-}{-}seconds}\NormalTok{ 20            }\CommentTok{\# JSON lines for N seconds}
\ExtensionTok{npm}\NormalTok{ run pull }\AttributeTok{{-}{-}} \AttributeTok{{-}{-}topic}\NormalTok{ us{-}east/hello                         }\CommentTok{\# the single latest message}
\end{Highlighting}
\end{Shaded}

Good for cron jobs, CI checks, shell glue, and quick manual probes.

\subsubsection{\texorpdfstring{4b. MCP server
(\texttt{axona-mcp})}{4b. MCP server (axona-mcp)}}\label{b.-mcp-server-axona-mcp}

\texttt{src/mcp.js} is a \textbf{Model Context Protocol} server that
exposes Axona as \textbf{native agent tools}, so Claude Code (and any
MCP-speaking agent) gets first-class \texttt{axona\_*} tools instead of
shelling out. It speaks JSON-RPC over \textbf{stdio} (stdout = protocol,
stderr = logs).

\textbf{It holds ONE persistent peer} (\texttt{src/mcp-session.js}),
opened lazily on the first call and kept connected --- so the agent is a
continuous participant, not a series of one-shot connections. That peer
signs with a \textbf{durable identity} (both the node identity and the
author keypair persist to
\texttt{\textasciitilde{}/.axona/claude-mcp-identity.json}), so the
agent keeps the \textbf{same Author ID and nodeId across restarts} --- a
stable on-network identity, not a fresh persona each call.

{\def\LTcaptype{none} % do not increment counter
\begin{longtable}[]{@{}
  >{\raggedright\arraybackslash}p{(\linewidth - 4\tabcolsep) * \real{0.3333}}
  >{\raggedright\arraybackslash}p{(\linewidth - 4\tabcolsep) * \real{0.3333}}
  >{\raggedright\arraybackslash}p{(\linewidth - 4\tabcolsep) * \real{0.3333}}@{}}
\toprule\noalign{}
\begin{minipage}[b]{\linewidth}\raggedright
Tool
\end{minipage} & \begin{minipage}[b]{\linewidth}\raggedright
Args
\end{minipage} & \begin{minipage}[b]{\linewidth}\raggedright
Purpose
\end{minipage} \\
\midrule\noalign{}
\endhead
\bottomrule\noalign{}
\endlastfoot
\texttt{axona\_publish} & \texttt{topic}, \texttt{message},
\texttt{region?} & Publish, signed by the stable author \\
\texttt{axona\_pull} & \texttt{topic}, \texttt{region?} & Fetch only the
latest message \\
\texttt{axona\_watch} & \texttt{topic}, \texttt{region?},
\texttt{since?}
(\texttt{all}\textbar{}\texttt{latest}\textbar{}\texttt{live}) & Open a
\textbf{standing subscription}; arrivals buffer server-side \\
\texttt{axona\_poll} & \texttt{topic?}, \texttt{peek?}, \texttt{max?},
\texttt{wait?}, \texttt{timeoutSec?} & Drain the buffer (the agent's
``inbox''); \texttt{wait:true} \textbf{long-polls} until a message
lands \\
\texttt{axona\_unwatch} & \texttt{topic}, \texttt{region?} & Stop a
watch, discard its buffer \\
\texttt{axona\_host} & \texttt{topic}, \texttt{region?} & \textbf{Root}
the topic on the agent's peer (store + serve, no subscribe) \\
\texttt{axona\_unhost} & \texttt{topic}, \texttt{region?} & Stop
hosting \\
\texttt{axona\_status} & --- & Peer nodeId + Author ID, mesh health,
per-watch + hosted lists \\
\texttt{axona\_subscribe} & \texttt{topic}, \texttt{region?},
\texttt{seconds?} (1--120), \texttt{since?} & One-shot listen window
(back-compat) \\
\end{longtable}
}

Region defaults to \texttt{useast} (\texttt{0x89}); subscribers must use
the \textbf{same region} as the publisher. The server's code default
targets \textbf{production}; set \texttt{RELAY\_NETWORK=testnet} to join
the staging network instead. Both run the 4.x line, but they are
separate networks --- the agent participates in whichever one its peer
joined, alongside the live apps on that network.

\textbf{Receiving --- push vs.~poll.} Two paths cover the same arrivals:
- \textbf{Push} --- every message on a watched topic is emitted as an
MCP \emph{logging notification} (\texttt{sendLoggingMessage}) to the
client. This is best-effort: it reaches the MCP client, but a turn-based
agent loop won't be interrupted by it. - \textbf{Long-poll} ---
\texttt{axona\_poll(wait:true,\ timeoutSec)} blocks \textbf{server-side}
until a message arrives (or the timeout), returning at push latency.
This is the reliable pull side. The standing watch buffers anything that
lands between polls, so nothing is lost in the gaps.

The agent becomes a \textbf{live participant} by pairing
\texttt{axona\_watch} with a \texttt{/loop} that calls
\texttt{axona\_poll(wait:true)} each iteration: watch once, then
poll-act-repeat. A message that asks for work is a request to execute
and reply; \textbf{approval for sensitive or irreversible actions stays
on the authenticated session, never the pub/sub channel} (a signed
message proves the author key, but the topic is still an external data
channel --- see the security boundary below).

Register it as a project-scoped MCP server with a \texttt{.mcp.json} at
your repo root:

\begin{Shaded}
\begin{Highlighting}[]
\FunctionTok{\{}
  \DataTypeTok{"mcpServers"}\FunctionTok{:} \FunctionTok{\{}
    \DataTypeTok{"axona"}\FunctionTok{:} \FunctionTok{\{} \DataTypeTok{"command"}\FunctionTok{:} \StringTok{"node"}\FunctionTok{,} \DataTypeTok{"args"}\FunctionTok{:} \OtherTok{[}\StringTok{"/abs/path/to/axona{-}relay/src/mcp.js"}\OtherTok{]} \FunctionTok{\}}
  \FunctionTok{\}}
\FunctionTok{\}}
\end{Highlighting}
\end{Shaded}

Point it at testnet (or a specific bridge) with an \texttt{env} block
--- also \texttt{MCP\_REGION} (default \texttt{useast}),
\texttt{MCP\_AUTHOR\_PATH} (identity file), \texttt{MCP\_BUFFER\_CAP}
(per-watch ring, default 1000):

\begin{Shaded}
\begin{Highlighting}[]
\FunctionTok{\{}
  \DataTypeTok{"mcpServers"}\FunctionTok{:} \FunctionTok{\{}
    \DataTypeTok{"axona"}\FunctionTok{:} \FunctionTok{\{}
      \DataTypeTok{"command"}\FunctionTok{:} \StringTok{"node"}\FunctionTok{,}
      \DataTypeTok{"args"}\FunctionTok{:} \OtherTok{[}\StringTok{"/abs/path/to/axona{-}relay/src/mcp.js"}\OtherTok{]}\FunctionTok{,}
      \DataTypeTok{"env"}\FunctionTok{:} \FunctionTok{\{} \DataTypeTok{"RELAY\_NETWORK"}\FunctionTok{:} \StringTok{"testnet"} \FunctionTok{\}}
    \FunctionTok{\}}
  \FunctionTok{\}}
\FunctionTok{\}}
\end{Highlighting}
\end{Shaded}

\begin{quote}
\textbf{Security boundary.} The MCP peer is a full
publisher/subscriber/host with a durable identity --- treat it as a real
network participant. Requests and reversible work can flow over a topic,
but the agent must keep the approval gate for destructive, irreversible,
or outward-facing actions on its authenticated session, not on an Axona
message: a signed envelope authenticates \emph{who}, but an open-topic
message is still untrusted external input.
\end{quote}

\subsubsection{\texorpdfstring{4c. Desktop app
(\texttt{axona-relay/desktop})}{4c. Desktop app (axona-relay/desktop)}}\label{c.-desktop-app-axona-relaydesktop}

An \textbf{Electron} GUI wrapper around the relay --- the same hosting +
metrics node, with a window instead of a terminal. Useful for
non-operators who want to contribute a durable node by leaving an app
open. Because it embeds a relay, a running desktop app adds one node to
your local fleet (and respawns it on crash).

\begin{center}\rule{0.5\linewidth}{0.5pt}\end{center}

\subsection{5. Directory \& federation}\label{directory-federation}

A single bridge is a single point of failure. Axona avoids that with two
cooperating behaviors, both built into the bridge's embedded peer:

\begin{itemize}
\tightlist
\item
  \textbf{Directory.} Each bridge with \texttt{BRIDGE\_DIRECTORY=on}
  publishes a signed entry
  (\texttt{\{url,\ lat,\ lng,\ label,\ ver,\ ts\}}) to the well-known
  open topic \texttt{axona:bridge-directory} (pinned to the
  \texttt{useast} region so everyone derives the same Topic ID) and
  \texttt{host()}s that topic so the entry survives. At launch a client
  collects the directory, ranks candidates (configured roots →
  known-good by reputation → fresh from the directory), and
  \textbf{fails over} to a saved alternate if its primary is
  unreachable.
\item
  \textbf{Federation.} A bridge with the directory on also
  \textbf{uplinks} --- an outbound \texttt{webTransport} connection
  \emph{into} a sibling bridge --- so the two bridges share one
  connectome. A client connected to \textbf{either} bridge can then
  discover and reach peers on \textbf{both}. In the live network,
  \texttt{bridge-west} uplinks to \texttt{bridge.axona.net}; the east
  bridge stays uplink-less as the seed.
\end{itemize}

\begin{quote}
Verify federation against a \textbf{warm} topic (both bridges host +
republish \texttt{axona:bridge-directory}), never a cold fresh topic ---
a cold topic has no roots yet and reads as a false zero.
\end{quote}

\begin{center}\rule{0.5\linewidth}{0.5pt}\end{center}

\subsection{6. The PoW collector}\label{the-pow-collector}

\texttt{pow-collector.js} (\texttt{npm\ run\ collect}) is a small
standing service that subscribes to the proof-of-work benchmark
\textbf{results} topic and archives each published result for later
analysis. It's optional infrastructure for the PoW research line, not
part of the core pub/sub path. (It can wedge silently --- alive but not
receiving --- so a restart with
\texttt{since:\textquotesingle{}all\textquotesingle{}} backfills the
retained backlog.)

\begin{center}\rule{0.5\linewidth}{0.5pt}\end{center}

\subsection{7. Choosing what to run}\label{choosing-what-to-run}

{\def\LTcaptype{none} % do not increment counter
\begin{longtable}[]{@{}
  >{\raggedright\arraybackslash}p{(\linewidth - 2\tabcolsep) * \real{0.5000}}
  >{\raggedright\arraybackslash}p{(\linewidth - 2\tabcolsep) * \real{0.5000}}@{}}
\toprule\noalign{}
\begin{minipage}[b]{\linewidth}\raggedright
You want to\ldots{}
\end{minipage} & \begin{minipage}[b]{\linewidth}\raggedright
Run
\end{minipage} \\
\midrule\noalign{}
\endhead
\bottomrule\noalign{}
\endlastfoot
Let strangers connect to your app's network & A \textbf{bridge} (or use
a live one --- \texttt{bridge.axona.net} for production,
\texttt{testnet.axona.net} for staging) \\
Keep a topic's messages alive when no author is online & A
\textbf{relay} that \texttt{host()}s it \\
Publish live subscriber counts for a topic & Nothing extra ---
\textbf{any root serves metrics on demand} when a client subscribes to
\texttt{metricTopic(T)}; a relay just makes the root durable \\
Publish/subscribe from a shell, cron, or CI & The \textbf{CLI}
(\texttt{axona-cli}) \\
Give an AI agent native Axona tools & The \textbf{MCP server}
(\texttt{axona-mcp}) \\
Contribute a node from a desktop, no terminal & The \textbf{desktop
app} \\
Survive a bridge outage & Two \textbf{federated bridges} + the
directory \\
\end{longtable}
}

Most application developers run \textbf{none} of these in development
--- they point the library at the live \texttt{bridge.axona.net} and let
the existing relay fleet host their topics. You reach for self-hosted
services when you want an isolated network, guaranteed topic durability,
or an agent/automation integration.

\begin{center}\rule{0.5\linewidth}{0.5pt}\end{center}

\subsection{8. Are there analogous
systems?}\label{are-there-analogous-systems}

If you know other realtime/decentralized stacks, the Axona services map
onto patterns you've already seen:

{\def\LTcaptype{none} % do not increment counter
\begin{longtable}[]{@{}
  >{\raggedright\arraybackslash}p{(\linewidth - 2\tabcolsep) * \real{0.5000}}
  >{\raggedright\arraybackslash}p{(\linewidth - 2\tabcolsep) * \real{0.5000}}@{}}
\toprule\noalign{}
\begin{minipage}[b]{\linewidth}\raggedright
Axona service
\end{minipage} & \begin{minipage}[b]{\linewidth}\raggedright
Closest analogue elsewhere
\end{minipage} \\
\midrule\noalign{}
\endhead
\bottomrule\noalign{}
\endlastfoot
Bridge (signaling + TURN) & WebRTC \textbf{STUN/TURN} servers;
\textbf{libp2p} circuit-relay v2 / rendezvous \\
Relay (durable topic host) & An \textbf{MQTT broker} / \textbf{NATS}
server with retained messages --- but decentralized and unprivileged \\
Bridge directory & \textbf{DNS seeds / bootstrap node lists} (Bitcoin,
IPFS) \\
MCP server & Any other \textbf{MCP integration} (GitHub, Slack) --- a
thin agent-facing adapter over the library \\
\end{longtable}
}

The difference in every row is the same: in those systems the server is
an \emph{authority} (it owns the namespace, the messages, or the
membership). In Axona the service is a \emph{volunteer} --- a peer that
took on a role and can be replaced by any other peer that takes on the
same role.

\begin{center}\rule{0.5\linewidth}{0.5pt}\end{center}

\subsection{Where to go next}\label{where-to-go-next}

\begin{itemize}
\tightlist
\item
  \textbf{\href{Axona-Programmer-Guide-v4.22.0.md}{Programmer Guide}}
  --- build the peer that talks to these services.
\item
  \textbf{\href{Axona-API-Reference-v4.22.0.md}{API Reference}} ---
  \texttt{host()} / \texttt{unhost()}, \texttt{metricTopic()}, and the
  rest of the public surface.
\item
  \textbf{\texttt{axona-bridge/deploy/INSTALL.md}} --- provision a
  bridge (installer, Docker, or manual).
\item
  \textbf{\texttt{axona-relay/README.md}} --- the relay's full
  configuration + dashboard reference.
\end{itemize}

Found a gap in this guide? Open an issue at
\url{https://github.com/axona-net/axona-docs}.

\end{document}
